\section{Teorema da Dualidade Árvore-Emaranhado}

Vamos agora ao nosso segundo teorema principal. Antes, algumas definições.

\begin{definicao*}[Estrela e estrela-vértice]
  $E \subset \mathcal{P}(V)$ é uma estrela se para quaisquer dois $A,B \in S$ temos $\lnot A \subset B$.

  Veja que $E$ é estrela se e somente se seu fecho pra complementar é $S(T)$ para algum $T$ que é uma estrela no sentido de grafos. (finita ou infinita!)

  Dado um conjunto árvore finito $S$, sabemos que $S \cong S(T)$ para uma única árvore $T$. Cada vértice $v$ de $T$ induz uma estrela $E_v$ em $S(T)$. Formalmente, $E_v$ é o conjunto de conjuntos $A$ de $S$ tal que a única aresta que liga $A$ a $\lnot A$ é uma aresta adjacente a $v$, e além disso $v \in A$. A imagem de $E_v$ em $S$ pelo isomorfismo é dito uma estrela-vértice do conjunto árvore $S$.
\end{definicao*}

\begin{proposicao*}[Definição intrínseca de estrelas-vértice]
  Tome $S(T)$ com $T$ finito ou infinito, e $E \subset S(T)$ uma estrela. Então, $E = E_v$ para um vértice $v$ de $T$ se e somente se ela satisfaz a `condição de minimalidade': não existe $A \in S(T) \setminus E$ e $X \subset E$ próprio ($X \neq E$) tais que $A \subset B$ pra todo $B \in X$ e $(E \setminus X) \cup \{A\}$ é uma estrela.
\end{proposicao*}
\begin{prova}
Suponha $E = E_v$. Tome um $A \in S(T) \setminus E$ qualquer, e um conjunto $X \subset E$ tal que $A \subset B$ pra todo $B \in X$. $A$ corresponde a alguma aresta $e$ da árvore $T$, necessariamente uma que não é adjacente a $v$ por $A \notin E$. Seja $f$ a aresta adjacente a $v$ que inicia o caminho do vértice $v$ até $e$. Seja $B$ o conjunto em $E$ que contém $v$ e está ligado a seu complementar via $f$. Então, $\lnot B$ contém ambos os vértices adjacentes a aresta $e$, ou seja, contém um elemento de $A$ e de $\lnot A$ necessariamente. Disto segue dois fatos: o primeiro é que é impossível $A \subset B$, e portanto $B \in E \setminus X$. O segundo é que é impossível $\lnot A \subset B$, e portanto $E \setminus X \cup \{A\}$ não pode ser uma estrela. 

Vamos provar a contrapositiva agora. Suponha que $E$ satisfaz a condição de minimalidade. Tome dois $A,B \in E$ quaisquer e sejam $e,f$ suas arestas em $T$ correspondentes. Seja $v$ a ponta da aresta $e$ em $A$ e $w$ a ponta da aresta $f$ em $B$. Queremos mostrar que $v = w$. Se eles fossem diferentes, então existe um caminho entre estes dois pontos. Seja $g$ a primeira aresta saindo de $v$ neste caminho. Seja $C$ a componente conexa de $T \setminus g$ que contém $v$. Veja que qualquer $D \in E$ será tal que ou $C \subset D$ ou $\lnot C \subset D$. Assim, se tomarmos como $X$ o conjunto de $D \in E$ tais que $C \subset D$, então teremos que $E \setminus X \cup \{C\}$ é uma estrela. Isto é proibido pela nossa hipótese inicial, e assim podemos concluir que $v = w$. Ou seja, todas as arestas vão ter o mesmo ponto inicial, digamos $v$. Pra finalizar a prova, basta vermos que todas as arestas adjacentes a $v$ tem uma de suas componentes conexas induzidas em $E$ - o que irá acontecer, pois se pra algum $e$ isto não acontecesse, bastaria adicionar a $E$ a componente conexa contendo $v$ e ainda teríamos uma estrela, contradizendo a condição de minimalidade (neste caso, $X = \varnothing$). Assim, $E = E_v$ como gostaríamos.

\end{prova}

\begin{pergunta}
Quando $T$ for uma árvore de ordem, a definição intrínseca de estrelas-vértice ainda funciona pra capturar apenas os $E_v$? Os vértices limite precisam ser considerados como tendo 'raios' adjacentes a eles?
\end{pergunta}

\begin{pergunta}
 A definição intrínseca de estrela-vértice no caso infinito ajuda a generalizar o teorema de dualidade pro caso infinito?
 \end{pergunta}

\begin{definicao*}[Evitar]
  Dizemos que um emaranhado $\tau$ de $S$ evita um conjunto $\mathcal{F}$ de subconjuntos de $S$ se nenhum $E \in \mathcal{F}$ é tal que $\tau$ seleciona todos os elementos de $E$.
\end{definicao*}


\begin{definicao*}[As hipóteses da dualidade]
  Dizemos que $\mathcal{F}$ decide $A \in S$ se ou $\{A\} \in \mathcal{F}$ ou $\{\lnot A\} \in \mathcal{F}$

  Um conjunto de estrelas $\mathcal{F}$ é dito $S$-consistente se sempre que $\{A\} \in \mathcal{F}$ então nenhum $B \subset A$, $B \in S$ é tal que $\{\lnot B\} \in \mathcal{F}$.
  
  Dado dois conjuntos $R \subset A$ definimos um mapa $\mathcal{P}(V) \rightarrow \mathcal{P}(V)$ que chamamos de $(R,A)$-relativização. Dado um $X$, o mapa leva

  \begin{align*} 
X &\mapsto  X \cap R \text{ se } X \subset A \\ 
X &\mapsto \lnot(\lnot X \cap R) \text{ se } \lnot X \subset A \\
X &\mapsto X \text{ caso contrário}
\end{align*}

  Dizemos que $S \subset \mathcal{P}(V)$ satisfaz a hipótese da dualidade se preserva $(R,A)$ relativizações pra quaisquer par $R,A \in S$ com $R \subset A$.
  
  Dizemos que um conjunto de estrelas $\mathcal{F}$ satisfaz a $S$-hipótese da dualidade se para quaisquer dois $A,B \in S$ com interseção vazia e não decididos por $\mathcal{F}$, existe um $R \in S$ tal que $R \subset A$, $\lnot R \subset B$, e a imagem de qualquer estrela em $\mathcal{F}$ via $(R,A)$-relativização e também via $(\lnot R,B)$-relativização é ainda uma estrela em $\mathcal{F}$.
\end{definicao*}


\begin{teorema*}[Dualidade Árvore-Emaranhados]
  Seja $S \subset \mathcal{P}(V)$ fechado para complementar, finito, e satisfazendo a hipótese da dualidade. Seja $\mathcal{F}$ um conjunto de estrelas de $S$ que seja $S$-consistente e satisfazendo a $S$-hipótese da dualidade.

  Então, vale apenas uma dentre as seguintes afirmações:

\begin{itemize}[label=\ding{212}]
  \item Existe um emaranhado de $S$ que evita $\mathcal{F}$
  \item Existe um conjunto árvore $S' \subset S$ cujas estrelas-vértices estão todas em $\mathcal{F}$
\end{itemize}
  \end{teorema*}
  \begin{prova}
    Primeiro provamos que as afirmações não botem ser ambas verdades. Isto acontece devido a um simples fato: dado uma árvore, qualquer escolha de orientação para todas as suas arestas (colocar uma seta na aresta em uma direção ou outra) é tal que existe pelo menos um vértice onde todas as setas apontam para ele. Retraduzindo este resultado grafista para nossa linguagem abstrata:

    \begin{lema*}
      Dada um emaranhado $\tau$ de um conjunto árvore $S(T)$, existe uma estrela-vértice cujos elementos são todos escolhidos por $\tau$.
    \end{lema*}
    \begin{prova}
      Um emaranhado induz, para cada aresta de $T$, uma orientação. Assim, agora temos uma árvore orientada. Chame de ralo um vértice cujas setas adjacentes todas apontam para ele. Começando a partir de qualquer vértice, se ele não for um ralo, existe uma aresta apontando para fora. Trace um caminho seguindo por arestas direcionadas para fora sempre que possível. Como a árvore é finita, este caminho deve terminar em algum vértice, e como árvores não possuem ciclos, este vértice final deve necessariamente ser um ralo. A estrela vértice referente a este ralo é tal que todos os seus elementos são escolhidos por $\tau$
    \end{prova}

    Suponha que tivessemos um emaranhado  $\tau$ de $S$ que evita $\mathcal{F}$ e também um conjunto árvore $S' \subset S$ com estrelas-vértices todas em $\mathcal{F}$. $\tau$ induz, por restrição, um emaranhado $\tau'$ de $S'$ que pelo lema anterior, deve ter um estrela-vértice cujos elementos são todos escolhidos por $\tau'$ e portanto $\tau$. Mas este estrela-vértice está em $\mathcal{F}$, contradizendo $\tau$ evitar $\mathcal{F}$.

    Agora seguimos para demonstrar que pelo menos uma dessas opções acontece. Faremos isto por indução no número de elementos de $S$ não decididos por $\mathcal{F}$. Suponha que não existe conjunto não decidido. Então podemos formar o conjunto $\{A \mid \{\lnot A\} \in \mathcal{F}\}$ como o único candidato possível para ser um emaranhado que evita $\mathcal{F}$. A consistência de $\mathcal{F}$ garante que ele é realmente um emaranhado, porém ele pode não evitar $\mathcal{F}$. Se este for o caso, existe uma estrela $E \subset S$ em $\mathcal{F}$ tal que todo $A \in E$ é tal que $\{\lnot A\} \in \mathcal{F}$. Como $E$ é estrela, $E$ também é conjunto árvore, referente a uma árvore que é uma estrela no sentido clássico de teoria de grafos - e todas as suas estrelas-vértices vão estar em $\mathcal{F}$ pois estas são apenas o próprio conjunto $E$ e os $\{\lnot A\}$. Então o teorema vale para $0$ conjuntos não decididos.
    
    Suponha que valha para $n$ conjuntos não decididos, e vamos supor que existem $n+1$ conjuntos não decididos. Tome um $A \in S$ não decidido. A ideia básica da demonstração é a seguinte: estenda $\mathcal{F}$ para $\mathcal{F}_1 \defeq \mathcal{F} \cup \{A\}$ e também para $\mathcal{F}_2 \defeq \mathcal{F} \cup \{\lnot A\}$. Se ambos estes conjuntos ainda forem consistentes e satisfazerem a hipótese da dualidade, então por eles conterem um conjunto decidido a menos, podemos aplicar o teorema. Um emaranhado que evita $\mathcal{F}_1$ ou $\mathcal{F}_2$ com certeza já evita $\mathcal{F}$. Agora, se estivermos no caso onde existem conjuntos árvore $S_1, S_2$ com suas estrelas em $\mathcal{F}_1$ e $\mathcal{F}_2$, iremos colar ambas estas árvores para formar uma única árvore com estrelas em $\mathcal{F}$, finalizando a demonstração. Vamos aos detalhes disso - esta construção será na verdade mais sutil do que a descrição que acabamos de dar.

    Para podermos garantir que $\mathcal{F}_1$ e $\mathcal{F}_2$ são consistentes e satisfazem a hipótese da dualidade, iremos ter que defini-los diferente do que afirmamos no parágrafo anterior. Seja $B_1$ maximal em $\{B \in S \mid A \subset B \text{ e } B \text{ não é decidido por }\mathcal{F}\}$, e $B_2$ maximal em $\{B \in S \mid \lnot A \subset B \text{ e } B \text{ não é decidido por }\mathcal{F}\}$. Defina então $\mathcal{F}_1 \defeq \mathcal{F} \cup \{\lnot B_1\}$ e $\mathcal{F}_2 \defeq \mathcal{F} \cup \{\lnot B_2\}$. 

    $\mathcal{F}_1$ é consistente: basta verificar que nenhum $Z \subset \lnot B_1$ é tal que $\{Z\} \in \mathcal{F}_1$ - o que é válido pois a existência de um $Z$ desse implicaria que $\lnot Z$ contradiz a maximalidade de $B_1$. Da mesma forma, $\mathcal{F}_2$ é consistente.

    $\mathcal{F}_1$ satisfaz a hipótese da dualidade: dados dois $A',B'$ de interseção não vazia e não decididos por $\mathcal{F}$, precisamos apenas garantir que a $(R,A')$-relativização (e a $(R,B')$ relativização) da estrela $\{\lnot B_1\}$ ainda é um elemento de  $\mathcal{F}_1$. Se a relativização preserva o $\lnot B_1$, então não há nada o que fazer. A maximalidade de $B_1$ garante que não podemos ter $B_1 \subset \lnot A'$. Resta apenas o caso em que $\lnot B_1 \subset \lnot A'$. Neste caso, a relativização da estrela resultará no conjunto $\{ \lnot B_1 \cap R \}$. É claro que $B_1 \subset \lnot( \lnot B_1 \cap R ) = B_1 \cup \lnot R$. A maximalidade de $B_1$ garante que $B_1 \cup \lnot R$ precisa ser decidido por $\mathcal{F}$. A consistência de $\mathcal{F}$ garante que $\{B_1 \cup \lnot R\} \notin \mathcal{F}$. Logo $\{\lnot B_1 \cap R\} \in \mathcal{F}$. É claro que, analogamente, provamos que a $(R,B')$ relativização também é preservada.

    Podemos então aplicar a hipótese de indução em $\mathcal{F}_1$, obtendo ou um emaranhado que evita $\mathcal{F}_1$, resolvendo nosso problema, ou obtendo um conjunto-árvore $S_1$ com todas as suas estrelas em $\mathcal{F}_1$. Se essas estrelas estiverem em $\mathcal{F}$, resolvido. Caso contrário, uma das estrelas será o $\{\lnot{B_1}\}$, e todas as outras em $\mathcal{F}$. 
    
    Analogamente, repetimos o raciocíno dos parágrafos anteriores agora com $\mathcal{F}_2$, e assim obtemos um conjunto-árvore $S_2$ com uma estrela sendo $\{\lnot{B_2}\}$ e todas as suas outras estrelas em $\mathcal{F}$. Visualmente, podemos imaginar duas árvores, uma com uma folha referente a $B_1$ e outra com uma folha refrente a $B_2$. Se estes conjuntos fossem complementar um do outro, poderamos colar estas duas árvores através de suas folhas e obtermos uma árvore com todas as estrels em $\mathcal{F}$ (como faremos ao final dessa demonstração). Como eles não necessariamente são complementares, vamos ter que modificar estas árvores para corrigir esse problema - é aqui que entra, finalmente, o uso das relativizações.

    Temos que $\lnot B_1$ e $\lnot B_2$ tem interseção vazia, e não são decididos por $\mathcal{F}$. Como $\mathcal{F}$ satisfaz a hipótese da dualidade, existe um $R \in S$ tal que $R \subset \lnot B_1$, $\lnot R \subset \lnot B_2$, e $\mathcal{F}$ preserva $(R,\lnot B_1)$-relativização e $(\lnot R, \lnot B_2)$-relativização.

    Aplique $(R,\lnot B_1)$-relativização em todos os elementos de $S_1$, obtendo $S_1'$. É uma conta de muitos casos, porém direta, verificar que $S_1'$ ainda se mantém árvore. Como $\mathcal{F}$ preserva estrelas após a relativização, $S_1'$ é um conjunto árvore cujas estrelas-vértices estão todas em $\mathcal{F}$, menos uma única estrela, que agora é referente ao conjunto $\{R\}$. Similarmente, ao aplicarmos $(\lnot R,\lnot B_2)$-relativização em $S_2$, obtemos um conjunto árvore $S_2'$ com todas as estrelas em $\mathcal{F}$ menos uma única estrela referente ao conjunto $\{\lnot R\}$.

    Agora, finalmente, podemos combinar $S_1'$ e $S_2'$ em um conjunto árvore $S'$ com todas as estrelas em $\mathcal{F}$. Olhe para as árvores $T_1$ e $T_2$ tais que $S_1' \cong S(T_1)$ e $S_2' \cong S(T_2)$. $T_1$ deve conter um vértice cuja estrela-vértice correspondente em $S_1'$ será $\{R\}$, e portanto este vértice deve necessariamente ser um vértice de grau $1$, o chamemos de $v_1$. Veja que toda aresta de $T_1$ separa a árvore de forma que uma das suas componentes conexas contém $\{v_1\}$. Como $\{v_1\}$ é a estrela que é mandada pra $\{R\}$ via o isomorfismo entre $S_1'$ e $S(T_1)$, concluimos que pra qualquer $X \in S_1$ teremos ou $R \subset X$ ou $R \subset \lnot X$. Analogamente, repetimos o mesmo raciocínio para $S_2'$ e sua árvore $T_2$, com seu vértice $v_2$ referente a estrela-vértice $\{\lnot R\}$. Daí, para qualquer $Y \in S_2'$, teremos ou $\lnot R \subset Y$ ou $\lnot R \subset \lnot Y$. Juntando estes dois fatos, concluimos que pra quaisquer $X \in S_1', Y \in S_2'$, estes dois conjuntos vão estar encaixados. Logo $S' \defeq S_1' \cup S_2'$ é um conjunto árvore. Sua árvore correspondente será a árvore $T'$ obtida a partir de $T_1$ e $T_2$ unindo ambas as árvores, as colando através dos vértices $v_1$ e $v_2$, e aí eliminando estes vértices, deixando sobrar uma aresta ligando diretamente o vizinho de $v_1$ em $T_1$ ao vizinho de $v_2$ em $T_2$. Assim, as estrelas-vértices de $T'$ são ou estrelas-vértices de $T_1$ fora $\{R\}$ ou estrelas vértices de $T_2$ fora $\{\lnot R\}$ - todas em $\mathcal{F}$, finalizando nossa demonstração.  

  \end{prova}

  Esse resultado é essencialmente um resultado sobre $\mathcal{F}$: sob a condição dele ser `consistente' e tendo `suficientes interseções', a única obstrução para juntarmos as estrelas em uma estrutura de árvore é a existência de uma orientação consistente dos subconjuntos adjacentes que compõe as estrelas de $\mathcal{F}$. Podemos refrasear o resultado sem menção a $S$, de duas formas diferentes, uma tomando o `menor' $S$ possível, e outra o `maior' $S$ possível.

\begin{teorema*}[Dualidade Árvore-Emaranhados  - minimal]
  Seja $\mathcal{F}$ um conjunto de estrelas de $\mathcal{P}(V)$, com $V$ finito. Seja $S$ o fecho pra complementar de $\cup \mathcal{F}$. Suponha que $\mathcal{F}$ seja $S$-consistente e satisfaça a $S$-hipótese da dualidade.

  Então, vale apenas uma dentre as seguintes afirmações:

\begin{itemize}[label=\ding{212}]
  \item Existe um emaranhado de $S$ que evita $\mathcal{F}$
  \item Existe um conjunto árvore $S' \subset S$ cujas estrelas-vértices estão todas em $\mathcal{F}$
\end{itemize}
  \end{teorema*}
  \begin{prova}
    Basta usar o teorema anterior com o $S$ sendo o fecho pra complementar de $\cup \mathcal{F}$ - como $\mathcal{F}$ satisfaz a $S$-hipótese da dualidade, o $S$ irá satisfazer a hipótese da dualidade.
  \end{prova}

  Ou então

  \begin{teorema*}[Dualidade Árvore-Emaranhados - maximal]
  Seja $\mathcal{F}$ um conjunto de estrelas de $\mathcal{P}(V)$, com $V$ finito. Suponha que $\mathcal{F}$ seja $\mathcal{P}(V)$-consistente e satisfaça a $\mathcal{P}(V)$-hipótese da dualidade.

  Então, vale apenas uma dentre as seguintes afirmações:

\begin{itemize}[label=\ding{212}]
  \item Existe um $v \in V$ que não pertence a $\cap E$ pra todo $E \in \mathcal{F}$
  \item Existe um conjunto árvore $S' \subset \mathcal{P}(V)$ cujas estrelas-vértices estão todas em $\mathcal{F}$
\end{itemize}
  \end{teorema*}
  \begin{prova}
    Basta usar o teorema com $S = \mathcal{P}(V)$ e lembrar que emaranhados sobre $\mathcal{P}(V)$ são todos ultrafiltros, e como esta é uma álgebra de Boole finita, são todos ultrafiltros principais.
  \end{prova}

  \begin{pergunta}
    O que acontece quando $S$ é infinito? Contraexemplos? Casos especiais onde o resultado vale? (localmente finito, rayless) ``Quais são as outras obstruções para formarmos uma árvore infinita a partir de infinitas estrelas?''
  \end{pergunta}

  \begin{pergunta}
    A versão maximal do teorema tem algum análogo em teoria de conjuntos clássica?
  \end{pergunta}

    \begin{pergunta}
   A apresentação de $S$ com todos os emaranhados sendo principais simplifica a demonstração ou a intuição por trás do teorema?
  \end{pergunta}
  


% Aqui pra baixo tem minha tentativa de pensar dualidade fraca pro infinito. Do jeito que tá escrito tem erros. O que eu sei agora: eu consigo provar que se existe o conjunto árvore, então toda orientação seleciona uma estrela ou uma extremidade. Também consigo provar que se toda orientação seleciona uma estrela, então existe um conjunto árvore. Mas não dá pra juntar esses dois num resultado só de dualidade bonitinha pro infinito.

  % \begin{definicao*}[Extremidades]
  %   Seja $S \subset \mathcal{P}(V)$ fechado pra complementar, $\epsilon$ um conjunto de elementos de $S$. Dizemos que $\epsilon$ é uma extremidade se $\epsilon = \{A_1, A_2, \dots\}$ e $A_{i+1} \subset A_i$ pra todo $i$.


  %   Dizemos que $\epsilon$ é uma extremidade de um conjunto de estrelas $\mathcal{F}$ se pra todo $i$ existe uma estrela $E$ em $\mathcal{F}$ tal que $\{A_i, \lnot A_{i+1}\} \subset E$. 
  % \end{definicao*}


  %   A volta vai ser a mesma demonstração do Jay, e notando que a orientação construida por elu evita extremidades. O jeito que eu fiz a demonstração delu foi construir um grafo a partir das estrelas de $\mathcal{F}$ (aresta entre estrela se dá pra juntar elas em uma árvore), mostrar que este grafo é uma floresta, descrever a condição neste grafo que garante a existência de um conjunto árvore com todas as estrelas-vértices em $\mathcal{F}$ e aí mostrar que se essa condição não vale podemos construir uma orientação que evita $\mathcal{F}$.



  % \begin{teorema*}[Dualidade fraca pro infinito]
  %   Seja $\mathcal{F}$ um conjunto de estrelas de um sistema $S \subset \mathcal{P}(V)$.  Então:

  %   Existe um conjunto-árvore $S(T) \cong S' \subset S$ cujas estrelas-vértices estão todas em $\mathcal{F}$ se e somente se toda orientação de $S$ seleciona uma estrela ou uma extremidade de $\mathcal{F}$.
  % \end{teorema*}
  % \begin{prova}
  %   Escrever a ida: essa é a versão infinita da demonstração que ambos os casos da dualidade não podem acontecer ao mesmo tempo (orientação de arestas em uma árvore ou cai num ralo ou vai pro infinito).

  %   Vamos para a volta. Primeiro resumimos a intuição da prova.Queremos construir um conjunto árvore a partir das estrelas de $\mathcal{F}$. Duas estrelas $E$ e $E'$ são `compatíveis', no sentido de que conseguimos colá-las em um processo de construção de uma árvore, apenas quando uma estrela possui um $A \in E$ e a outra o complementar $\lnot A \in E$. Este processo de colagem pode ser feito até chegarmos em uma árvore cujas estrelas são todas de $\mathcal[F]$. Se em toda possível tentativa desta construção falhar, por falta de estrelas em $\mathcal{F}$, isto será suficiente para construirmos uma orientação de $S$ que evita $\mathcal{F}$.

  %   Para formalizarmos esta construção, definimos um grafo cujos vértices são as estrelas de $\mathcal{F}$. Duas estrelas $E$ e $E'$ serão adjacentes se existe um $A \in S$ tal que $A \in E$ e $\lnot A \in E'$ - e colorimos esta aresta com uma cor associada ao par $\{A,\lnot A\}$. As cores $\{A, \lnot A\}$ pra cada $A \in E$, são ditas as cores de $E$. Se $E$ tem pelo menos uma aresta pra cada uma de suas cores, dizemos que $E$ é uma estrela feliz - caso contrário, ela é uma estrela triste. 
    
  %   Primeiramente, note que este grafo é uma árvore.

  %   Uma subárvore feliz $T_{f}$ é uma subárvore onde pra cada vértice $E$ dela: o $E$ é uma estrela feliz, e $T_f$ contém uma e apenas uma aresta pra toda cor de $E$. Se convença disso: achar uma subárvore feliz é a mesma coisa que encontrar um conjunto árvore $S' \subset S$ cujos estrela-vértice estão todos em $\mathcal{F}$.

  %   Suponha que não exista subárvore feliz - vamos conseguir construir uma orientação $\tau$ de $S$ que não seleciona estrela nem extremidade de $\mathcal{F}$. Vamos olhar pra cada $A \in S$ e decidir se $\tau$ escolhe ele ou seu complementar. Pros $A \in S$ tais que $A$ nem $\lnot A$ estão em alguma estrela de $\mathcal{F}$, $\tau$ pode escolher qualquer um dos $A, \lnot A$.  Resta decidirmos os conjuntos que estão em estrelas. Para isso, tome uma estrela $E$ de $\mathcal{F}$. Primeiramente suponha que ela é triste. Então existe um $A \in E$ tal que $\lnot A \notin E'$ pra todo $E' \in \mathcal{F}$. Decida então que $\tau(\lnot A) = 1$ - isto já evita que $\tau$ escolhe $E$, e não toma nenhuma decisão sobre os conjuntos das outras estrelas. Agora suponha que $E$ é feliz. Faça agora uma construção indutiva em cada componente conexa de $\mathcal{F}$.
  % \end{prova}



  %Aqui é eu de novo tentando formalizar o algoritmo de 'apontar pro problema', acho que tá certo mas foi antes de eu ver a demonstração mto mais simples e direta de Jay

% %\begin{proposicao*}[Dualidade Árvore-Emaranhado Fraca]
% Seja $\mathcal{B}$ uma álgebra de Boole e considere $\mathcal{E} \subset \mathcal{P}(\mathcal{B})$. Seja $S$ o fecho pra complementar de $\cup \mathcal{E}$, e suponha que ele seja finito. Chamamos os elementos de $\mathcal{E}$ de estrelas. Dizemos que duas estrelas $E_1,E_2 \in \mathcal{E}$ se colam se existe $a \in E_1, \lnot a \in E_2$. Suponha que não exista uma sequência $E_1, E_2, \dots, E_n$ tal que $E_i$ cola com $E_{i+1}$ e $E_n$ cola com $E_1$, pra $1 \leq n $. Então existe um subconjunto $T \subset \mathcal{E}$ tal que $\cup T$ é fechado pra complementar se e somente se toda orientação $\tau$ de $S$ satisfaz alguma estrela de $\mathcal{E}$.
% \end{proposicao*}

% \begin{prova}
% Considere $\mathcal{E}$ como um grafo com arestas coloridas, onde duas estrelas $E_1$ e $E_2$ são ligadas por uma aresta colorida por $\{a, \lnot a\} \subset S$ se elas se colam via $a$, ou seja, $a \in E_1$ e $\lnot a \in E_2$. Por simplicidade, vamos fazer referência a cor $\{a, \lnot a\}$ usando apenas um dos dois elementos, digamos $a$. A hipótese de não existência de ciclos garante que $\mathcal{E}$ é uma floresta.

% Chame uma estrela $E$ de triste se existe um $a \in E$ tal que $\lnot a \notin E'$ pra todo $E' \in \mathcal{E}$. Uma estrela é feliz se ela não é triste. Seja $C_E$ o conjunto de cores de arestas adjacentes a uma estrela $E$. Veja que $T \subset \mathcal{E}$ é tal que $\cup T$ é fechado pra complementar se e somente se $T$ é uma subfloresta onde, pra toda estrela $E \in T$, $E$ é feliz e $T$ contém arestas adjacentes a $E$ de todas as cores em $C_E$. Podemos tomar um $T' \subset T$ que é uma árvore e contém, pra cada $E \in T'$ exatamente uma aresta de cada cor de $C_E$. Vamos chamar um $T$ dessa forma de solução.

% ($\implies$) Suponha que exista um $T \subset \mathcal{E}$ com $\cup T$ fechado pra complementar, e seja $T' \subset T$ como descrito no parágrafo anterior. Uma orientação $\tau$ de $S$ induz uma escolha de direção pra cada aresta de $\mathcal{E}$: digamos que se $\tau(a) = 1$, então uma aresta colorida por $\{a,\lnot a\}$ aponta na direção da estrela que contém $a$. Assim, $\tau$ induz uma escolha de direção pra cada aresta de $T'$. Como $T'$ é uma árvore finita, se escolhermos um vértice qualquer dela e traçarmos um caminho direcionado maximal, chegaremos a um vértice onde todas as suas arestas adjacentes apontam para dentro - este vértice é, então, exatamente uma estrela satisfeita por $\tau$.

% ($\impliedby$) Suponha que não exista $T \subset \mathcal{E}$ com $\cup T$ fechado pra complementar. Vamos construir uma orientação $\tau$ de $S$ que não satisfaz nenhuma estrela. Se uma estrela $E$ é triste, então temos um $a \in E$ com $\lnot a \notin E'$ pra todo $E' \in \mathcal{E}$ - defina $\tau(\lnot a) = 1$, e assim não satisfazemos $E$ e como $\lnot a \notin E'$, não `atrapalhamos' a não satisfazibilidade de nenhuma outra estrela. Definido $\tau$ para todas as estrelas tristes, agora começamos a defini-lo para as felizes. Faremos isso através de um processo indutivo. Tome uma estrela $E_0$ feliz. Deve existir um $a \in E_0$ tal que pra qualquer outra estrela $E'$ com $\lnot a \in E'$, temos que $E' \setminus {\lnot a}$ não é extendível através de uniões com estrelas de $\mathcal{E}$ para um conjunto fechado pra complementar. Caso contrário iriamos unir todas estas estrelas, referentes aos $E'$ de cada $a \in E_0$, junto com $E_0$, e formaríamos um $T$ fechado pra complementar. Defina $\tau( \lnot a) = 1$. Para toda estrela $\tilde{E}$ que admite um caminho até $E_0$ que não passa por uma aresta colorida por $a$, seja $\{b,\lnot b\}$ a cor da aresta saindo de $\tilde{E}$ iniciando o caminho, e se $b \in \tilde{E}$, oriente $\tau (\lnot b) = 1 $. A não existência de ciclos garante que não há conflitos nestas definições de $\tau$ - ela está bem definida até agora. Seja $\mathcal{E'}$ o conjunto de estrelas que admitem um caminho para $E_0$ passando por uma aresta colorida por $a$, porém substitua uma estrela $E'$ que contém $\lnot a$ pela estrela $E' \setminus {\lnot a}$. Agora, este conjunto não pode admitir um $T \subset \mathcal{E'}$ com $\cup T$ fechado pra complementar (se tivesse, $T$ teria que ter alguma das estrelas modificadas $E' \setminus {\lnot a}$, mas já supomos que nenhuma delas é extendível para algum conjunto fechado pra complementar). Por indução, podemos então garantir a existência de um $\tau'$ que orienta o fecho pra complementar de $\cup \mathcal{E'}$ e não satisfaz nenhuma estrela de $\mathcal{E'}$. Assim, extendemos $\tau$ com as orientações vindas de $\tau'$ (novamente sem conflitos devido a não existência de ciclos). Faça este mesmo processo para outras componentes conexas de $\mathcal{E}$, e assim temos um $\tau$ bem definido que não satisfaz nenhuma estrela.